\documentclass[../../../include/open-logic-section]{subfiles}

\begin{document}

\olfileid[fr]{sth}{replacement}{limofsize}
\olsection{Limitation de taille}

La tentative peut-être la plus courante de justification
«~intrinsèque~» du remplacement passe par le principe suivant :
\begin{enumerate}
	\item[] \limofsize. Des objets quelconques forment un ensemble,
	pourvu qu'ils ne soient pas trop nombreux.
\end{enumerate}
Sous l'axiome de choix, ce principe justifie immédiatement le remplacement. Un ensemble formé par remplacement ne peut
alors pas être plus grand que l'ensemble à partir duquel il a été formé. Plus précisément, supposons
que le remplacement permette de former l'ensemble $\funimage{\tau}{A} = \Setabs{\tau(x)}{x \in A}$ ; sous cette hypothèse, $\cardle{\funimage{\tau}{A}}{A}$.\footnote{Note
de l'édition française : la comparaison source est définie par l'existence d'une injection. Sans hypothèse supplémentaire, la
surjection de $A$ sur son image ne suffit pas en général à fournir une injection en sens
inverse ; le choix suffit ici. Pour des contre-exemples dans $\ZF$, voir G. Shen, \emph{Cantor's theorem may
fail for finitary partitions} (2023), théorème 3.3 et corollaire 4.17, \url{https://arxiv.org/abs/2309.00235}.} Si les !!{element}s de $A$ n'étaient
pas trop nombreux pour former un ensemble, leurs images ne sont alors pas trop nombreuses pour former $\funimage{\tau}{A}$.

La difficulté évidente, lorsqu'on invoque \limofsize{} pour
justifier le remplacement, est que nous n'avons \emph{pas}
encore posé de principe tel que \limofsize. De plus, dans
notre récit de la conception cumulative et itérative des
ensembles, présenté dans les \crefrange{sth:story::chap}{sth:z::chap},
rien ne \emph{suggérait} \limofsize. C'est précisément
pourquoi Boolos a écrit un jour : «~Peut-être peut-on conclure
qu'il y a au moins deux idées “derrière” la théorie des
ensembles~» \citeyearpar[p.~19]{Boolos1989}. D'un côté, les
idées qui sous-tendent la conception cumulative et itérative
des ensembles sont censées justifier~$\Z$. De l'autre,
\limofsize{} est censée justifier le remplacement.

La difficulté n'est cependant pas seulement que nous sommes
restés jusqu'ici \emph{silencieux} sur \limofsize. C'est aussi
que \limofsize{} (sous la forme que nous venons de lui donner)
semble assez mal s'accorder avec la conception cumulative et
itérative des ensembles. Après tout, ce principe ne dit rien
de la formation des ensembles par \emph{étapes}.

Ce décalage n'a rien de très surprenant, étant donné l'histoire
de ces «~deux idées~» (la conception cumulative et itérative des
ensembles et \limofsize). Elles correspondent, en fin de compte,
à deux projets assez différents pour écarter les paradoxes de
la théorie des ensembles. La conception cumulative et itérative
écarte le paradoxe de Russell en disant, en gros :
\emph{nous n'aurions jamais dû nous attendre à ce qu'un ensemble
de Russell existe, puisqu'il ne serait «~formé~» à aucune étape}.
En revanche, \limofsize{} entend exclure cet ensemble en disant,
en gros : \emph{nous n'aurions jamais dû nous attendre à ce qu'un
ensemble de Russell existe, puisqu'il serait trop grand}.

Présentons les choses sans détour : envisagée comme une réponse
aux paradoxes, \limofsize{} demande \emph{beaucoup} plus de
justification. Considérons, par exemple, cette version du
paradoxe de Russell : \emph{aucun carlin ne renifle exactement
les carlins qui ne se reniflent pas eux-mêmes} (voir la
\olref[sth][story][rus]{sec}). Si vous demandez «~pourquoi
n'existe-t-il pas un tel carlin ?~», répondre qu'il lui
faudrait renifler trop de carlins n'est pas une bonne réponse.
Pourquoi serait-ce alors une bonne explication intuitive
de la non-existence d'un ensemble de Russell que de dire
qu'il serait «~trop grand~» pour exister ?

En somme, on peut vous pardonner de rester quelque peu perplexe
devant la motivation «~intuitive~» de \limofsize.

\end{document}
